\documentclass[11pt,a4paper]{article}
\usepackage[margin=25mm]{geometry}
\usepackage{amsmath,amssymb,amsthm,hyperref}
\setlength{\emergencystretch}{2em}
\newtheorem{theorem}{Theorem}
\newtheorem{corollary}[theorem]{Corollary}
\newcommand{\gen}[1]{\langle#1\rangle}
\title{Completion of the ambient paired-head branch}
\author{Human Compression project: L044 continuation}
\date{}
\begin{document}\maketitle
Retain the complete standing hypotheses and smaller-degree induction
of \emph{Critical odd-degree tails behind exceptional ends}.
Thus $G\le A_n$ has a connected string C-marking, $n\ge12$,
and every maximal ambient omission is intransitive. Eligibility below
always refers to that ambient omission, not to a restricted tail.

\begin{theorem}
In either of the following source-head configurations, suppose that
each colour $j\ge k$ has two ambient-eligible edges within some
single actual $H=G_{\ge k}$-orbit:
\begin{enumerate}
\item $G_0$ has the two orbits $\{a\}$ and its complement;
$k=2$ and $b=a^{\rho_0}$.
\item The first three generators are the actual supported six-point
head from source Table 2(2) or (3), in an extendable orientation,
with $k=5$. The accepted exceptional-head interface is retained:
$D=\{1,\ldots,6\}$, $\tau=\rho_3=(6\ b)\delta$,
$H$ fixes $D\cup\{b\}$, and $G_{\ge4}$ is transitive on
$\Omega\setminus D$.
\end{enumerate}
Then $r\le\lfloor(n-1)/2\rfloor$.
\end{theorem}
\begin{proof}
Suppose the rank bound fails. The complete critical-head reduction
gives $n=2r$, with the actual natural alternating tail used by the
exceptional-head proofs. In the singleton case $n\ge14$, and
$H=G_{\ge2}$ fixes $a,b$ and is intransitive on its support.
In the exceptional case $n\ge18$ and $H$ fixes its seven-point head.

Choose one such orbit and pair for each tail colour. The whole-orbit
transport theorem normalizes this assignment so that its full local
$K=G_{>k}$ is intransitive in each assigned orbit. Every colour is
assigned exactly once; the assigned colour sets are disjoint.
For the actual nontrivial $H$-orbits $O_s$ put
$d_s=|F_s|-|O_s|/2$, and let $f=|\operatorname{Fix}(H)|$.
The normalized bounds give all $d_s\le0$ and
\[
                    S=\sum_s d_s=f/2-k\ge-3/2.       \tag{1}
\]

Let $\tau=\rho_0$ in the singleton case and $\tau=\rho_3$
in the exceptional case. This involution centralizes $H$ and
belongs to every ambient omission of a tail colour.
On the specified fixed head it is the single transposition
$(a\ b)$ or $(6\ b)$, respectively. Since $\tau$ is even,
its action on the complement of that head is odd.

Centralization permutes actual $H$-orbits equivariantly.
An exchanged pair of nontrivial orbits contributes an odd sign
only when their common size is odd; the accepted equivariant
aggregation theorem then gives $S\le-3$, contrary to (1).
On a preserved nontrivial orbit, a nonidentity centralizing
involution is free. Its sign is odd only on an orbit of size
$2q$ with $q$ odd. If the assignment is nonempty, the complete
strict-internal and localized-internal proofs exclude it in both
heads, including the final first-assigned odd-anchor configuration.
If the assignment is empty and $q\ge3$, its exact deficit
is $-q\le-3$, again contrary to (1). The case $q=1$ is the
empty internal pair excluded by the localized two-point theorem.

It remains to include the $H$-fixed points, which are absent
from its support. The full ambient interval theorem gives
\[
            2(r-k)\le n-f-c,\qquad f+c\le2k,          \tag{2}
\]
where $c\ge1$ counts the nontrivial $H$-orbits.
In the singleton case $f\ge2$ and $c\ge2$ force $f=c=2$;
there are no additional fixed points outside $\{a,b\}$.
In the exceptional case, if the action on the additional fixed
points is odd, it exchanges at least two of them. The complete
localized two-point theorem excludes exactly this possibility;
its reduction from (2) gives $f=9,c=1$ before applying the actual
pure-three-cycle/IP obstruction.

These cases exhaust every odd contribution to $\tau$ outside
its head: preserved nontrivial orbits, exchanged nontrivial orbits,
and exchanged singleton orbits. Each is impossible. This contradicts
the required odd sign, proving the bound.
\end{proof}

\begin{corollary}[The precise residual source interface]
Under either critical head and rank failure, some tail colour has
at most one ambient-eligible edge in every actual $H$-orbit.
\end{corollary}
\begin{proof}
Otherwise choose an orbit with two eligible edges independently
for each colour and apply the theorem. Such a choice does not
require different colours to use different orbits.
\end{proof}

\paragraph{Attribution and limits.}
This combines the complete paid orbit-transfer, equivariant-deficit,
strict-internal and localized-interval proofs for the all-paired
part of Cameron--Fernandes--Leemans--Mixer,
\emph{Highest rank of a polytope for $A_n$}, arXiv:1605.09173v1,
Proposition 5.31, physical pages 42--44. It uses the source's
deficit strategy, with the stated replacement proofs and an explicit
accounting of singleton orbits. No stronger literal local-deficit
or diagram assertion is inferred from the final contradiction.

The condition is that each colour has an eligible pair \emph{in one
actual orbit}. Two edges split across different orbits do not
satisfy it. Nor does a pair crossing only a local tail omission.
The corollary identifies, rather than resolves, that remaining
source-routing branch. Complete source recovery, rigidity/table
contracts, formalization and original or question credit remain
outside this note.
\end{document}
